\lhead{\emph{Conclusion Générale}}
\chapter*{Conclusion Générale}

L'objectif principal de ce projet était de développer une solution décisionnelle intelligente capable d'exploiter les données de l'entreprise BIOSPHERE afin d'améliorer le suivi des performances et d'optimiser la gestion des stocks grâce à des mécanismes d'analyse et de prévision.

Les travaux réalisés ont permis de mettre en place une chaîne décisionnelle complète couvrant l'ensemble du cycle de valorisation des données. Une architecture de collecte et d'intégration a été développée afin d'alimenter un Data Warehouse centralisant les informations stratégiques de l'entreprise. Cette infrastructure a constitué la base nécessaire à la production d'indicateurs fiables et à la réalisation des analyses décisionnelles.

L'exploitation de ces données a conduit à la création de tableaux de bord permettant aux responsables de disposer d'une vision synthétique de l'activité à travers différents indicateurs liés aux ventes, aux achats, aux marges, aux profits et aux niveaux de stock. Ces outils facilitent le suivi des performances et contribuent à améliorer la qualité des décisions opérationnelles.

Une composante prédictive a également été intégrée au système à travers le développement de modèles de Machine Learning. Les modèles de classification ont permis d'identifier les produits présentant un risque de rupture de stock tandis que les modèles de régression ont été utilisés pour estimer les niveaux futurs de stock et générer des recommandations de réapprovisionnement. Ces mécanismes offrent une approche proactive de la gestion des stocks en permettant d'anticiper les situations critiques avant leur apparition.

L'ensemble des fonctionnalités développées a ensuite été intégré dans une application desktop décisionnelle regroupant les modules de Business Intelligence, de prévision des stocks, de gestion des alertes et d'administration. Cette application fournit aux utilisateurs une plateforme unique d'aide à la décision adaptée aux besoins de l'entreprise.

Au-delà des résultats techniques obtenus, ce projet a permis de démontrer la complémentarité entre les approches de Business Intelligence et de Machine Learning. L'association de l'analyse descriptive et de l'analyse prédictive constitue un levier important pour améliorer le pilotage des activités et renforcer les capacités d'anticipation des décideurs.

Néanmoins, certaines limites subsistent. La solution développée repose principalement sur les données historiques disponibles et demeure déployée sous forme d'application desktop. De futures évolutions pourraient porter sur le développement d'une version web, l'intégration de nouvelles sources de données externes, l'automatisation du réentraînement des modèles prédictifs ou encore l'utilisation d'algorithmes plus avancés basés sur le Deep Learning.

En définitive, ce projet a permis de répondre aux objectifs fixés en proposant une plateforme décisionnelle intelligente capable de transformer les données de BIOSPHERE en informations stratégiques exploitables, contribuant ainsi à une meilleure maîtrise des stocks et à une prise de décision plus efficace.
